\label{ch:impl}

\section{Development Environment}
Training was done in Python with PyTorch, scikit-learn and XGBoost. The API dependency file is shorter: FastAPI, Uvicorn, a CPU build of PyTorch and TorchVision, pandas, requests and Pydantic. The client is React, bundled with Vite, with Leaflet for the map and Axios for the API. Photographs are stored with the browser's IndexedDB. Versions are pinned for the torch packages in the deploy file so the hosted build does not drift.

Locally, the API is started with \path{uvicorn api.main:app --reload --port 8000} after the Groq key is set in the shell. The client reads \path{VITE_API_BASE} and otherwise uses \path{http://localhost:8000}. A separate flag turns off the in-memory photo-check cap for a local evaluation only. It is not set on the host.

\section{Data and Training Programs}
\texttt{scripts/collect\_soil\_data.py} and \texttt{scripts/collect\_weather\_data.py} write the raw point tables. \texttt{scripts/synthesize\_dataset.py} joins them to crops and applies the rule engine. \texttt{scripts/validate\_dataset.py} checks the joined table. \texttt{scripts/prepare\_final\_dataset.py} writes the train, validation and test splits and the scaler. \texttt{scripts/train\_baseline\_models.py} fits Random Forest and XGBoost. \texttt{scripts/train\_nn.py} fits AgroDynamicNet, including the branch ablations. Repeated seeds write a report per seed and do not overwrite the weights the API loads.

\texttt{scripts/train\_wheat\_classifier.py} fine-tunes either ResNet18 or MobileNetV3-Small. The 40 pilot files are excluded by hash. The remaining images are split once, with a fixed split seed, into 671 training, 144 validation and 144 test images, and that split is reused for every model seed. The script records accuracy, macro-F1, a temperature, an expected calibration error before and after scaling, and a coverage table at several probability thresholds. \texttt{scripts/gradcam\_wheat.py} can draw a Grad-CAM map for the ResNet18 checkpoint \cite{selvaraju2017gradcam}. That script is an analysis tool. The website does not call it.

\texttt{scripts/sample\_eval.py} draws the 40-image pilot, ten images each of yellow rust, brown rust, powdery mildew and healthy leaves. Septoria is one of the five training classes and is not in the pilot folders. \texttt{scripts/eval\_diagnose.py} sends the pilot through the photo-check endpoint once, without any change to the prompt in response to the scores.

\section{How the Two Models Are Fit}
The depth labels are built before any network is trained. Soil and weather are joined on location and season. Each crop is attached only to its season. The rule of Chapter~\ref{ch:design} produces a depth, noise of standard deviation 0.15\,cm is added, and the row is clipped. The training script then standardises the numeric inputs with means and standard deviations stored beside the split, so the API can apply the same scale at inference. Crop name and state are integer codes into the two embedding tables. A branch that the ablation turns off is replaced by a zero vector of width 32 after that branch would have run, which keeps the trunk the same shape for every variant.

The wheat classifier starts from an ImageNet checkpoint. Training images are randomly cropped to $224 \times 224$ with a scale down to half the frame, flipped on either axis, rotated by up to 20 degrees, and jittered in brightness, contrast, saturation and hue. Validation, test and pilot images are only resized to $224 \times 224$ and normalised with the ImageNet mean and standard deviation. The 40 pilot files are removed by an MD5 hash of the file bytes before the split, so a renamed copy does not leak back into training. The remaining images are split once, stratified by class, with split seed 42: 70\% train and 30\% held out, and that held-out half is split again into validation and test. The counts are 671, 144 and 144. Every model seed reuses the split. What changes with the seed is the initialisation noise, the order of the training batches, and the dropout masks.

The classifier loss is cross-entropy with class weights inverse to the training counts, and with label smoothing of 0.05. The optimiser is AdamW. The pretrained backbone uses learning rate $10^{-4}$ and the new classification head uses $10^{-3}$, with weight decay $10^{-4}$. A cosine schedule runs over a maximum of 20 epochs. The checkpoint that is restored is the one with the best validation macro-F1, and training stops after six epochs without an improvement. Temperature scaling is fit afterwards, on the validation logits of that checkpoint, by minimising cross-entropy over a single positive temperature with LBFGS. The temperature does not change which class wins. It changes the probability that is compared with 0.8. Expected calibration error on the test split uses ten equal-width bins of that probability. The API loads the seed-0 ResNet18 weights and the temperature stored with them, 0.8484. MobileNetV3-Small is trained in the same harness so that the five-seed table has a smaller baseline. Its weights are not loaded.

Random Forest and XGBoost are fit on the same depth split as comparators. The forest has 300 trees with a maximum depth of 12. The boosted model has 300 trees, a maximum depth of 6 and a learning rate of 0.05. Both use a fixed random seed. A five-fold cross-validation MAE is computed on the training rows only, as a check that the test number is not a single lucky partition. Those models are not served. Serving them would have given up the branch flags, which are the reason the ablation and the deployed network are the same architecture.

\section{API Modules}
\path{api/main.py} loads AgroDynamicNet. It exposes \path{/health}, \path{/crops}, \path{/predict} and \path{/predict_by_location}, and includes the other routers. \texttt{/predict} accepts soil and weather fields directly, which is how a developer can bypass the live lookups. \texttt{/predict\_by\_location} is what the page uses.

\texttt{api/retrieval.py} tokenises notes, expands the Hindi and Hinglish map, and ranks with BM25. \texttt{api/chat.py} applies the rate limit, builds the prompt and returns the answer with the notes that were used. \texttt{api/diagnose.py} loads the problem list, calls the vision model, filters identifiers, and, when the crop is wheat, calls \texttt{api/wheat\_classifier.py}. That module decodes the image, applies the ImageNet transform, scales the logits by the stored temperature, and sets \texttt{abstained} when the top probability is below 0.8. \texttt{api/forecast.py} is the server-side forecast used in local tests. The page no longer calls it.

\section{Client Modules}
The depth screen is a map and a form. Choosing a point stores the latitude and longitude in the browser so that My crop can show the same forecast. The result card prints the network depth, the rule depth, the safe range, the texture, the estimated moisture, and a note when cached soil was used. The chat panel sits under that card and sends the latest result as context. Suggested questions are starting points only.

My crop is a second view, not a second site. Before a plan exists it asks for the crop and the date. After that it lists tasks in date order with a status chip and a checkbox. The forecast card is the first block on that view as well. The photo timeline accepts a camera capture or a file, compresses it, and shows the check button only for wheat. The result component renders the PAU matches and, when present, the ResNet18 block with its probability or its withheld message.

The advisor layout places the map and the forecast in the left column and the form, the depth card and the chat in the right column. On a narrow screen the same blocks stack in that order. This is a presentation change. It does not change a prediction.

\section{One Depth Request, End to End}
A click on the map stores the latitude and longitude in the page and in browser storage. The form sends the crop, the state, the season and that point to \path{/predict_by_location}. The API checks that the season is known and that the point lies inside the bounding box of the selected state. It then asks SoilGrids for sand, clay, pH, organic carbon, bulk density and nitrogen, with a short timeout and a single retry. If any of those fields is missing, it replaces the soil record with the nearest cached point and keeps the distance.

Climatology for the same point is read next. The season's row supplies temperature, humidity, rainfall, the rain flag, solar radiation, wind and the sowing month. Phosphorus, potassium and moisture are estimated from the soil record by the same functions used in training. Those values, together with the crop's seed size and seed weight, are scaled with the stored training means and standard deviations and passed through AgroDynamicNet. The rule engine is called on the unscaled values. The response contains both depths, the safe range, the texture class, the moisture estimate, and a note. The note is where a cached-soil distance is written. The page does not choose between the two depths. It prints them side by side.

A photo check is a different request. The browser resizes the image and sends it, with the crop name, to \path{/diagnose}. For wheat, the server runs the vision-model match against the 14 identifiers and, separately, \path{classify_base64} on the ResNet18. The JSON that comes back has the PAU candidates and a classifier object with the label, the probability, the threshold, and whether the call was withheld. The page renders those fields. It does not recompute the threshold.

\section{What the Page Does When a Service Fails}
The free hosts fail in ways a demonstration has to absorb. SoilGrids is queried with a short timeout and one retry. If sand, clay, pH, organic carbon, bulk density or nitrogen is missing, the API replaces the soil record with the nearest cached point and writes the distance into the note. If the cache itself is empty, the request fails with a sentence that says live soil data is unavailable. Weather for the depth model is the stored climatology for that point and season. If that lookup fails, the page says the weather could not be fetched, and it does not invent a climate in the browser.

A point outside the rectangle of the selected state is rejected before either lookup. The message names the state and asks for a click inside it, or for the dropdown to be changed. Chat without a key on the server says that chat is not configured. A photo check without a key says the same about the photo check. If the language or vision service is unreachable, or returns nothing usable, the page says so and asks for another try. It does not fill the gap with a depth-only answer pretending to be a diagnosis.

The photo-check counter allows 6 requests from one address in a window of 600 seconds, and 40 in a day. Chat allows 10 in the same window and 80 in a day. Both counters are in the process memory. Restarting the API clears them, which is useful on a developer machine and is also why the provider's own quota is the limit that matters on the host. A local evaluation of the 40-image pilot sets a flag that disables the photo-check counter. That flag is not set in the deployed environment.

The forecast route on the API is still present and returns a clear failure when the host cannot reach Open-Meteo. The live page does not call that route. The browser calls Open-Meteo itself, which succeeds from the networks on which the host's outbound call did not. If the browser call fails, the card says the forecast is unavailable and leaves the depth card as it was. The next-task line is computed from the plan in local storage, so it can still name a task when the weather request fails, and it stays silent when no plan has been saved.

\section{What Is Stored Where}
Table~\ref{tab:storage} is the privacy boundary of the system. It is also the reason a plan does not follow the farmer to a second phone.

\begin{table}[H]
\centering
\caption{Where each piece of farmer data lives.}
\label{tab:storage}
\small
\begin{tabularx}{\textwidth}{lX}
\toprule
\textbf{Data} & \textbf{Where it is kept} \\
\midrule
Soil cache and model weights & On the API host, built before deployment \\
Chat and photo-check requests & Sent for an answer, not written to an application database \\
Crop plan & \texttt{localStorage} on the device \\
Field photographs & IndexedDB on the device, long side at most 1280\,px \\
Last map point & \texttt{localStorage}, so the forecast can reload \\
Groq credential & Host environment variable, not in the repository \\
\bottomrule
\end{tabularx}
\end{table}

\section{Repository Layout}
\texttt{scripts/} holds collection, synthesis, training and evaluation. \texttt{models/} holds the AgroDynamicNet definition and the two weight files the API loads. \texttt{api/} holds the HTTP layer. \texttt{frontend/} holds the client, including \texttt{src/data/schedules.json}. \texttt{data/knowledge/} holds the paraphrased notes. \texttt{data/problems/} holds the wheat problem list. \texttt{data/processed/} holds the metric reports for every seed. The wheat photographs themselves are not in the repository.
